\chapter{指数函数与对数函数}\label{ch:g12:exp}

指数函数是唯一与其自身导数相等的函数，
并以 $1$ 为其在 $0$ 处的值。它把和化为积；其逆函数
即自然对数，把积化为和。二者共同刻画
凡变化率与其大小成正比的现象：
放射性衰变、人口增长、复利。

\section{指数函数}

\begin{theorem}[存在性与唯一性]\label{thm:g12:exp:existence}
存在唯一的可导函数 $f \colon \R \to \R$，使得
\[
f' = f \qquad\text{且}\qquad f(0) = 1 .
\]
该函数称为\emph{指数函数}\index{指数函数}，记作
$\exp$，或 $x \mapsto \eu^x$。
\end{theorem}

\begin{proof}[唯一性的证明]
首先，这样的函数永不为零。事实上，设
$g(x) = f(x)f(-x)$，则
\[
g'(x) = f'(x)f(-x) - f(x)f'(-x) = f(x)f(-x) - f(x)f(-x) = 0,
\]
因此 $g$ 为常数，且 $g(0) = 1$：对每个 $x$，
$f(x)f(-x) = 1$，特别地 $f(x) \neq 0$。

现设 $f_1, f_2$ 为两个解，令 $h = \dfrac{f_1}{f_2}$（
由于 $f_2$ 永不为零）。于是
\[
h' = \frac{f_1' f_2 - f_1 f_2'}{f_2^2} = \frac{f_1 f_2 - f_1 f_2}{f_2^2} = 0,
\]
因此 $h$ 为常数，且 $h(0) = 1$，\emph{即} $f_1 = f_2$。

\emph{存在性}在本阶段不加证明地承认（它可以
由\cref{exo:g12:seq:10}中的数列得到，或作为
在\cref{ch:g12:integ}中通过积分构造的对数函数的反函数得到）。
\end{proof}

\begin{proposition}[函数方程]\label{prop:g12:exp:funceq}
对所有 $x, y \in \R$ 及 $n \in \Z$：
\[
\eu^{x+y} = \eu^x\,\eu^y, \qquad
\eu^{-x} = \frac{1}{\eu^x}, \qquad
\eu^{x-y} = \frac{\eu^x}{\eu^y}, \qquad
\bigl(\eu^{x}\bigr)^n = \eu^{nx}.
\]
此外，$\eu^x > 0$ 对所有 $x$ 成立。
\end{proposition}

\begin{proof}
固定 $y$，考虑 $\varphi(x) = \dfrac{\exp(x+y)}{\exp(x)\exp(y)}$。
其分子与分母作为 $x$ 的函数，都是
$f' = f$ 的解（在相差常数 $\exp(y)$ 的意义下）；直接对 $\varphi$ 求导
（利用商的求导法则）得 $\varphi' = 0$，因此
$\varphi \equiv \varphi(0) = 1$，即得第一个等式。取
$y = -x$ 即得第二式（其中 $\eu^0 = 1$），
第三式随之可得。第四式由第一式对 $n \geq 0$ 用归纳法得到，
再由第二式推广到 $n < 0$。

正性：$\eu^x = \left(\eu^{x/2}\right)^2 \geq 0$，且 $\eu^x \neq 0$
（见\cref{thm:g12:exp:existence}），故 $\eu^x > 0$。
\end{proof}

\begin{proposition}[变化与极限]\label{prop:g12:exp:variations}
指数函数严格递增、凸，且
\[
\lim_{x\to-\infty} \eu^x = 0, \qquad
\lim_{x\to+\infty} \eu^x = +\infty .
\]
\end{proposition}

\begin{proof}
由 $(\exp)' = \exp > 0$ 得严格递增；由 $(\exp)'' = \exp > 0$ 得
凸性。由凸性，$\eu^x \geq 1 + x$（在 $0$ 处的切线，见
\cref{ex:g12:deriv:ineq}），由比较可得 $\eu^x \to +\infty$（当 $x \to +\infty$ 时）。
再由 $\eu^{x} = 1/\eu^{-x} \to 0$（当 $x \to -\infty$ 时）。
\end{proof}

\begin{omfigure}
\begin{tikzpicture}
\begin{axis}[omaxis,
  width=8.8cm, height=6cm,
  xmin=-3.3, xmax=2.4, ymin=-2.4, ymax=7.8,
  xtick={-2,-1,1,2}, ytick={1}]
  \addplot[omProp, thick, domain=-3.2:2.3] {1 + x};
  \addplot[omDef, thick, domain=-3.2:2.0] {exp(x)};
  \fill (axis cs:0,1) circle (1.5pt);
  \node[omDef, font=\small, anchor=east] at (axis cs:1.35,5.2) {$\eu^x$};
  \node[omProp, font=\small, anchor=north west] at (axis cs:1.05,1.6) {$y = 1+x$};
\end{axis}
\end{tikzpicture}

{\small 指数函数：严格递增、凸，且
$\lim_{-\infty} \eu^x = 0$。其图像位于它在 $0$ 处的切线上方：
$\eu^x \geq 1 + x$。}
\end{omfigure}

\begin{theorem}[增长比较]\label{thm:g12:exp:growth}
\index{增长比较}
对每个整数 $n \geq 1$，
\[
\lim_{x\to+\infty} \frac{\eu^x}{x^n} = +\infty,
\qquad
\lim_{x\to-\infty} x^n \eu^x = 0 .
\]
换言之，指数函数的增长快于 $x$ 的任意次幂。
\end{theorem}

\begin{proof}
当 $n = 1$ 时，应用 $\eu^t \geq 1 + t \geq t$，并取 $t = x/2$，
\[
\frac{\eu^x}{x} = \frac{\left(\eu^{x/2}\right)^2}{x}
\geq \frac{(x/2)^2}{x} = \frac{x}{4} \xrightarrow[x\to+\infty]{} +\infty .
\]
对一般的 $n$，写作
\[
\frac{\eu^x}{x^n}
= \frac{1}{n^n}\left(\frac{\eu^{x/n}}{x/n}\right)^{n}
\xrightarrow[x\to+\infty]{} +\infty
\]
由 $n=1$ 的情形及复合即得上式。$-\infty$ 时的极限
由代换 $x \mapsto -x$ 得到：
$\abs{x^n \eu^x} = \frac{\abs{x}^n}{\eu^{\abs x}} \to 0$.
\end{proof}

\begin{omfigure}
\begin{tikzpicture}
\begin{axis}[omaxis,
  width=8.8cm, height=6cm,
  xmin=-0.4, xmax=5.6, ymin=-12, ymax=155,
  xtick={1,2,3,4,5}, ytick={50,100,150}]
  \addplot[omThm, thick, domain=0:5.4] {x^2};
  \addplot[omProp, thick, domain=0:5.3] {x^3};
  \addplot[omDef, thick, domain=-0.4:5.0] {exp(x)};
  \fill (axis cs:4.536,93.4) circle (1.5pt);
  \draw[black, dashed, thin] (axis cs:4.536,0) -- (axis cs:4.536,93.4);
\end{axis}
\end{tikzpicture}

{\small 指数函数（蓝色）最终超过每一个幂函数：它
在较早处超过 $x^2$（红色），超过 $x^3$（橙色）则发生于 $x \approx 4.54$ 处。}
\end{omfigure}

\section{自然对数}

\begin{definition}[自然对数]\label{def:g12:exp:ln}
指数函数连续且严格递增，是从 $\R$ 到
$\intoo{0}{+\infty}$ 的满射；由双射定理
（\cref{thm:g12:limcont:bijection}），对每个 $y > 0$，方程
$\eu^x = y$ 有唯一解。该解称为\emph{自然对数}\index{自然对数}，
$y$ 的，记作 $\ln y$。于是
\[
\text{对 } x \in \R,\ y > 0: \qquad y = \eu^x \iff x = \ln y .
\]
特别地，$\ln 1 = 0$，$\ln \eu = 1$，并且
$\eu^{\ln y} = y$, $\ln(\eu^x) = x$.
\end{definition}

\begin{omfigure}
\begin{tikzpicture}
\begin{axis}[omaxis,
  width=8.6cm, axis equal image,
  xmin=-4.3, xmax=6.3, ymin=-4.3, ymax=6.3,
  xtick={1}, ytick={1}]
  \addplot[black, dashed, domain=-3.9:5.9] {x};
  \addplot[omDef, thick, domain=-4.1:1.79] {exp(x)};
  \addplot[omThm, thick, domain=0.0167:5.9, samples=120] {ln(x)};
  \node[omDef, font=\small, anchor=east] at (axis cs:0.95,2.6) {$\eu^x$};
  \node[omThm, font=\small, anchor=north] at (axis cs:4.3,1.25) {$\ln x$};
  \node[font=\small, anchor=west, rotate=45] at (axis cs:4.6,4.85) {$y=x$};
\end{axis}
\end{tikzpicture}

{\small $\exp$ 与 $\ln$ 的图像关于
直线 $y = x$ 互为镜像：点 $(x, \eu^x)$ 反射为 $(\eu^x, x)$。}
\end{omfigure}

\begin{proposition}[代数性质]\label{prop:g12:exp:lnalg}
对所有 $a, b > 0$ 及 $n \in \Z$：
\[
\ln(ab) = \ln a + \ln b, \quad
\ln\frac{1}{a} = -\ln a, \quad
\ln\frac{a}{b} = \ln a - \ln b, \quad
\ln(a^n) = n \ln a, \quad
\ln\sqrt{a} = \tfrac12 \ln a .
\]
\end{proposition}

\begin{proof}
$\eu^{\ln a + \ln b} = \eu^{\ln a}\eu^{\ln b} = ab$，两边取 $\ln$
即得第一式。其余等式同理，由
\cref{prop:g12:exp:funceq}中的相应恒等式推出。
\end{proof}

\begin{proposition}[分析性质]\label{prop:g12:exp:lnanalytic}
函数 $\ln$ 在 $\intoo{0}{+\infty}$ 上可导，且
\[
(\ln x)' = \frac{1}{x},
\]
严格递增、凹，并且
$\lim\limits_{x\to0^+} \ln x = -\infty$,
$\lim\limits_{x\to+\infty} \ln x = +\infty$。此外，对每个整数
$n \geq 1$,
\[
\lim_{x\to+\infty} \frac{\ln x}{x^{1/n}} = 0, \qquad
\lim_{x\to0^+} x \ln x = 0 .
\]
\end{proposition}

\begin{proof}
反函数的可导性在本阶段不加证明地承认；承认此点后，
对恒等式 $\eu^{\ln x} = x$ 用链式法则求导：
$(\ln x)'\,\eu^{\ln x} = 1$，故 $(\ln x)' = \frac{1}{x} > 0$，从而严格
递增，又 $(\ln x)'' = -\frac1{x^2} < 0$，从而为凹。
$0^+$ 与 $+\infty$ 处的极限由双射与 $\exp$ 的极限对应得到。对于增长比较，
取 $x = \eu^{t}$：
$\frac{\ln x}{x} = \frac{t}{\eu^t} \to 0$，当 $t \to +\infty$ 时，由
\cref{thm:g12:exp:growth}，类似地，对另一极限
取 $x = \eu^{-t}$，得 $x\ln x = -t\eu^{-t} \to 0$。
\end{proof}

\begin{method}[用 $\exp$ 与 $\ln$ 解方程]\label{met:g12:exp:equations}
两个函数都严格递增，因此可以对方程或不等式两边同时作用
（或从两边去掉）它们，而不改变其方向：
\[
\eu^{u} = \eu^{v} \iff u = v, \qquad
\eu^{u} \leq \eu^{v} \iff u \leq v,
\]
对正数上的 $\ln$ 同样成立。\emph{解题前务必先检查
定义域}（$\ln$ 要求真数为正）。对形如
$a^x = b$（$a>0$，$a \neq 1$）的方程，改写为 $a^x = \eu^{x\ln a}$，
并解出 $x = \frac{\ln b}{\ln a}$。
\end{method}

\begin{example}\label{ex:g12:exp:halflife}
某放射性样本按 $N(t) = N_0\,\eu^{-\lambda t}$ 衰减，其
\emph{半衰期} $T$ 满足 $N(T) = N_0/2$，\emph{即}
$\eu^{-\lambda T} = \frac12$，从而 $T = \frac{\ln 2}{\lambda}$：半衰期
与初始量无关。
\end{example}

\section{练习}

\begin{exercise}[$\star$]\label{exo:g12:exp:1}
化简
$\dfrac{\eu^{3x}\,\eu^{-x+1}}{\eu^{x}}$ 与
$\ln\!\left(\dfrac{\eu^2\sqrt{\eu}}{\eu^{-3}}\right)$。
\end{exercise}

\begin{exercise}[$\star$]\label{exo:g12:exp:2}
在 $\R$ 中解下列方程：
\[
\text{(a) } \eu^{2x} - 3\eu^{x} + 2 = 0;
\qquad
\text{(b) } \ln(x - 1) + \ln(x + 2) = \ln 4 .
\]
\end{exercise}

\begin{exercise}[$\star$]\label{exo:g12:exp:3}
求下列极限：
\[
\lim_{x\to+\infty} \frac{\eu^x - x^2}{\eu^x + 1}, \qquad
\lim_{x\to+\infty} \frac{\ln x}{\sqrt x}, \qquad
\lim_{x\to0^+} x^2 \ln x, \qquad
\lim_{x\to+\infty} \bigl(x - \ln x\bigr).
\]
\end{exercise}

\begin{exercise}[$\star\star$]\label{exo:g12:exp:4}
研究函数 $f(x) = x\,\eu^{-x}$ 在 $\R$ 上的变化情况、极限、
极值；证明其曲线有一个拐点，并求出它的
坐标。
\end{exercise}

\begin{exercise}[$\star\star$]\label{exo:g12:exp:5}
证明：对所有 $x > -1$，$\ln(1 + x) \leq x$，且等号仅在
$x = 0$ 时成立。由此推出，对所有 $n \geq 1$，
\[
\left(1 + \frac{1}{n}\right)^n \leq \eu \leq \left(1 - \frac{1}{n+1}\right)^{-(n+1)} .
\]
\end{exercise}

\begin{exercise}[$\star\star$]\label{exo:g12:exp:6}
一笔本金 $C_0$ 以年利率 $3\%$ 投资，
利息每年按复利计算。
\begin{enumerate}
  \item 求本金 $C_n$ 在 $n$ 年后的表达式。
  \item 经过多少年本金翻一番？先用 $\ln$ 给出精确值，
        再给出数值。
  \item 与银行家所用的近似规则相比较，该规则为“用 $70$ 除以
        年利率（以百分数表示）”。
\end{enumerate}
\end{exercise}

\begin{exercise}[$\star\star$]\label{exo:g12:exp:7}
解不等式 $\eu^{2x} - \eu^{x+1} > 0$，再解不等式
$\ln(x^2 - 1) \leq \ln(x + 5)$。
\end{exercise}

\begin{exercise}[$\star\star\star$]\label{exo:g12:exp:8}
设 $f(x) = \dfrac{\eu^x}{x}$，其中 $x > 0$。
\begin{enumerate}
  \item 研究 $f$ 在 $\intoo{0}{+\infty}$ 上的变化情况，并求其
        最小值。
  \item 当 $k$ 取何值时，方程 $\eu^x = kx$ 有 $0$、$1$
        或 $2$ 个解（位于 $\intoo{0}{+\infty}$ 内）？
\end{enumerate}
\end{exercise}

\begin{exercise}[$\star\star\star$]\label{exo:g12:exp:9}
对 $n \geq 1$，设 $u_n = \left(1 + \frac1n\right)^n$。
\begin{enumerate}
  \item 利用 \cref{exo:g12:exp:5}，证明 $(u_n)$ 有
        上界 $\eu$。
  \item 证明 $\ln u_n = n \ln\left(1 + \frac1n\right) \to 1$，并
        由此推出 $u_n \to \eu$。
        （提示：识别 $\ln$ 在 $1$ 处的差商。）
\end{enumerate}
\end{exercise}

\section{问题：对数驯服世界}

\begin{problem}[{周末问题 --- 翻倍时间与
72 法则，地震、酸度与钢琴的标度，以及
常数 $\eu$ 处处留下痕迹}]
\label{pb:g12:exp:1}
凡是按固定\emph{百分比}增长的东西，都呈指数增长
--- 如存款、细菌、疫情 --- 而跨越过多数量级、
难以把握的量 --- 如地震能量、
酸度、声强 --- 都由对数来驯服。
本题计算翻倍时间，并揭示银行家的
72 法则，读取世界的对数标度，并收集
数 $\eu$ 在每一处犯罪现场留下的指纹
（见\cref{prop:g12:exp:funceq}、\cref{thm:g12:exp:growth}、
\cref{met:g12:exp:equations}）。

\medskip
\noindent\textbf{第一部分 --- 熟练运算。}
\begin{enumerate}
  \item 解方程：$\eu^{2x} = 5$；\quad $\ln(3x - 1) = 2$；\quad
        $\eu^{2x} - 3\eu^x + 2 = 0$（伪装的二次方程）。
  \item 化简：$\dfrac{\ln 8}{\ln 2}$；\quad
        $\ln\!\left(\eu^3 \sqrt{\eu}\right)$；\quad
        $\eu^{\ln 5 - \ln 2}$。
  \item 证明母不等式：$\eu^x \geq 1 + x$ 对所有
        实数 $x$ 成立（研究 $f(x) = \eu^x - 1 - x$），并由此推出
        $\ln(1 + u) \leq u$，其中 $u > -1$。
  \item 计算
        $\lim_{x \to +\infty} x^2 \eu^{-x}$，
        $\lim_{x \to +\infty} \frac{\ln x}{x}$
        （见\cref{thm:g12:exp:growth}），以及
        $\lim_{x \to 0} \frac{\eu^x - 1}{x}$（识别出
        一个导数）。
  \item 研究 $f(x) = x \ln x$ 在 $\intoo{0}{+\infty}$ 上的
        单调性、最小值，以及 $0^+$ 处的极限（在本阶段不加证明地承认：
        $x \ln x \to 0$）。这个小函数用来度量
        信息与熵，贯穿各本科分册。
\end{enumerate}

\medskip
\noindent\textbf{第二部分 --- 翻倍时间与
72 法则。}
\begin{enumerate}[resume]
  \item 存款每年增长 $3\,\%$。解方程
        $1.03^n = 2$：需要多长时间才能使本金翻倍？
  \item 银行家估计翻倍时间的公式为
        $\frac{72}{\text{利率 }\%}$。检验该法则在
        $3\,\%$、$6\,\%$ 与 $9\,\%$ 处的情形，并与精确值
        $\frac{\ln 2}{\ln(1 + r)}$ 比较，再加以解释：当
        $r$ 较小时，$\ln(1 + r) \approx r$（第 3 题的不等式只是
        故事的一半），因此精确常数为
        $100 \ln 2 \approx 69.3$ --- 银行家为何偏爱
        $72$？
  \item 一种细菌每 $20$ 分钟分裂一次：
        $p(t) = 2^{t/20}$（$t$ 以分钟为单位）。将其改写为
        $\eu^{\lambda t}$ 的形式，再计算 $p$ 在 $24$ 小时后的值。
        该答案（超过 $10^{21}$）说明了这一模型的何种问题
        --- 又有什么制约着真实的菌落？
  \item 咖啡因在体内的半衰期约为 $5$
        小时。于 15:00 饮用 $100$ mg 咖啡后，
        23:00 时还剩多少？它何时降至
        $10$ mg 以下？（失眠也有一个对数。）
  \item 一欧元按年利率 $100\,\%$ 计息：计算
        分别按年、按月及按日复利时的年末本金，
        并给出 \cref{exo:g12:exp:9} 所证明的
        不可突破的上限。哪个著名
        常数是纯粹贪婪的极限？
  \item 科学家为何要在疫情初期将 $\ln(\text{病例数})$ 关于
        时间作图？该图上的一条
        直线揭示了什么，
        其斜率度量什么？
\end{enumerate}

\medskip
\noindent\textbf{第三部分 --- 世界上的对数标度。}
（记 $\log_{10} x = \frac{\ln x}{\ln 10}$。）
\begin{enumerate}[resume]
  \item 以分贝计的声级：$L = 10 \log_{10}(I/I_0)$。一段对话测得
        $60$ dB，一场摇滚音乐会测得 $120$ dB，
        两者的声强相差多少倍？
  \item 地震震级每增加 $1$，
        振幅便变为原来的 $10$ 倍，释放的能量与振幅$^{3/2}$成正比。比较
        震级为 $5$ 与 $7$ 的两次地震：
        先比较振幅之比，
        再比较能量之比。
  \item 化学：$\text{pH} = -\log_{10}[\mathrm{H^+}]$。
        柠檬汁的 pH 为 $2$，牛奶的 pH 为 $6.5$：它们的
        酸浓度之比是多少？
  \item 音乐：每个八度使频率翻倍。从钢琴最低的
        A（$27.5$ Hz）到最高的 C
        （$4\,186$ Hz），键盘跨越了多少个八度？
        而我们的感官 --- 听觉、视觉与对地震的感受 ---
        为何偏爱对数标度？（用一句话回答。）
  \item 计算尺，即工程师使用了 $350$
        年的计算工具：两根刻度尺，其\emph{长度}
        （到刻度 $x$ 处）与 $\ln x$ 成正比。解释
        如何通过将一根尺沿另一根滑动来实现数的乘法，
        并指出起作用的恒等式，
        即 \cref{prop:g12:exp:lnalg}。
\end{enumerate}

\medskip
\noindent\textbf{第四部分 --- $\eu$ 的指纹。}
\begin{enumerate}[resume]
  \item 由 \cref{exo:g12:exp:8} 可知，方程
        $\eu^x = kx$（$x > 0$）的解的个数为 $0$、$1$ 或 $2$，
        取决于 $k$ 相对于某一
        阈值的位置。重述此结论，
        并验证其背后的几何事实：直线 $y = \eu x$
        恰是过原点的指数函数的
        切线。
  \item “险些中奖”常数：设有 $n$ 张彩票，
        每张的中奖概率均为 $\frac1n$，
        则 $\P(\text{未中奖}) = \left(1 - \frac1n\right)^n$。
        求其极限，借助
        $n \ln\!\left(1 - \frac1n\right)$（利用标准极限
        $\frac{\ln(1 + u)}{u} \to 1$），并与
        \cref{pb:g11:binom:1} 中神秘的 $0.368$ 相互印证。
  \item “$37\,\%$ 法则”：要从随机顺序面试的 $n$
        位候选人中选出最优者（不可回头），
        最优策略是拒绝前
        $\frac{n}{\eu} \approx 37\,\%$ 位，再选取
        第一位优于此前所有候选人者 --- 其成功
        概率约为 $\frac1\eu$。第 18 题与第 19 题中的两个
        $\frac1\eu$ 从何而来 ---
        请说明共同机制（许多独立的
        小概率事件），并将 $\frac1\eu$ 精确到
        小数点后三位。
  \item 终章 --- $\eu$ 的画像：复利的上限
        （第 10 题）；在 $0$ 处切线斜率
        恰为 $1$ 的底数（第 3 题与第 4 题）；没有
        多项式能追上的增长（第 4 题）；
        险些中奖与最优停止的常数（第 18--19 题）；
        以及它的逆 $\ln$，它把乘积变为和
        （第 16 题），把十倍的跨度变为英寸（第三部分）。
        每项各用一句话作答。
\end{enumerate}
\end{problem}
